(1) 记 \(f=1-3a/(2r)+a^3/(2r^3)\)、\(g=1-3a/(4r)-a^3/(4r^3)\)。散度为
\[
\nabla\!\cdot\mathbf u=U\cos\theta[f'+2(f-g)/r]=0.
\]
\(f(a)=g(a)=0\)，而 \(f(\infty)=g(\infty)=1\)，两项边界条件均满足。将速度代入向量Laplacian，\(-\nabla p+\mu\nabla^2\mathbf u=0\) 给出
\[
p(r,\theta)=p_\infty-\frac{3\mu Ua}{2r^2}\cos\theta.
\]
可对各径向幂次分别求导再代回方程独立核对。

(2) 于 \(r=a\) 有 \(\partial_ru_r=0\)，所以
\[
\sigma_{rr}=-p_\infty+\frac{3\mu U}{2a}\cos\theta,\qquad
\sigma_{r\theta}=\mu(\partial_ru_\theta+r^{-1}\partial_\theta u_r-u_\theta/r)
=-\frac{3\mu U}{2a}\sin\theta.
\]
取从球面指向流体的法向 \(\mathbf e_r\)，流体对球的轴向应力为
\(t_z=\sigma_{rr}\cos\theta-\sigma_{r\theta}\sin\theta\)。
去掉积分为零的均匀压强项，得到
\[
t_z=\frac{3\mu U}{2a},\quad
F_{\rm p}=2\pi\mu aU,\quad F_{\rm shear}=4\pi\mu aU,
\quad\boxed{F_z=6\pi\mu aU}.
\]
压力部分使用球面积 \(\int\cos^2\theta\,d\Omega=4\pi/3\)，剪切部分用 \(\int\sin^2\theta\,d\Omega=8\pi/3\)；流体力沿来流方向。

(3) Stokes方程的应力散度为零，点乘速度并用不可压性得到
\(\nabla\cdot(\sigma\mathbf u)=2\mu E:E\)。对球体外流体体积积分，由球面无滑移和远场机械功输入得
\[
\mathcal D=\int_{r>a}2\mu E:E\,dV=F_zU
=\boxed{6\pi\mu aU^2}.
\]
也可换到球相对无穷远静液以速度 \(-U\mathbf e_z\) 运动的Galilean系；此时阻力耗散功率为 \(F_{\rm drag}U\)。稳态动能不累积不等于没有供能：正定剪切耗散由外部维持来流的机械功持续补足。

(4) 终端运动时，净重力
\((4\pi a^3/3)(\rho_s-\rho_f)g\)
被 \(6\pi\mu aV_t\) 抵消，故对于重球
\[
V_t=\frac{2a^2(\rho_s-\rho_f)g}{9\mu}.
\]
以直径计的雷诺数
\(Re=2\rho_faV_t/\mu\ll1\) 是忽略惯性的必要自洽条件。若有显著壁面阻塞、湍流或球形状变化，则应修正该线性拖曳公式。
